\section{Specificaties}
\label{sec:specificaties}

Op basis van de probleemstelling en de wensen van MYLAPS zijn de specificaties 
en onderzoekscriteria voor het systeem vastgesteld. Het project 
richt zich op een softwarematige implementatie binnen de bestaande hardware van 
de ProRTK tracker. \\

Omdat de paardensport het primaire scenario is, zijn de dynamische eigenschappen van een galopperend paard een belangrijk punt. Volgens The Jockey Club is de hoogst gemeten snelheid in de paardensport tot wel 70,76 km/u~\cite{jockeyclub_racehorses}. Het systeem moet 
deze snelheden met een veiligheidsmarge kunnen verwerken, wat resulteert in 
een minimale snelheidseis van 80 km/u. Naast de hoge snelheden resulteren de dynamische bewegingen tot eisen aan het meetbereik van de sensoren. \\

Uit biomechanisch onderzoek~\cite{witte2006effect} blijkt dat bij 
galopperende racepaarden op snelheden van 9 tot 17 m/s een verticale kracht 
van 24,7 N/kg op de voorledematen wordt gemeten bij topsnelheid. Via 
$F = m \cdot a$ volgt hieruit een piekversnelling op het zwaartepunt van 
ongeveer 2,5 $g$. De versnelling die op de romp wordt gemeten ligt lager, omdat 
de romp als demper werkt ten opzichte van de ledematen. De voorwaartse en 
zijwaartse versnellingen liggen bovendien aanzienlijk lager dan de verticale 
piekwaarde. Het accelerometerbereik van de IMU moet minimaal ruim boven deze 
2,5 g liggen om clipping te voorkomen. \\

Het gyroscoopbereik is bepaald op basis van de verwachte hoeksnelheden tijdens gebruik. Bij een object dat met snelheid $v$ een cirkel met straal $r$ aflegt, geldt voor de hoeksnelheid formule \ref{eq:hoeksnelheid}.

\begin{equation}
    \omega = \frac{v}{r} \text{ [rad/s]}
    \label{eq:hoeksnelheid}
\end{equation}

Normaliter specificeren datasheets van gyroscopen de hoeksnelheid in graden per seconde (DPS, \textit{Degrees Per Second}), waardoor de uitkomst van formule \ref{eq:hoeksnelheid} moet worden omgerekend. Omdat een volledige cirkel $2\pi$ radialen of 360 graden bevat, wordt de conversie van rad/s naar DPS berekend via formule \ref{eq:conversie}:

\begin{equation}
    \text{DPS} = \omega \times \frac{180}{\pi}
    \label{eq:conversie}
\end{equation}

Voor de straal van de bochten is als eerste uitgegaan van de geometrie van een standaard racebaan in Amerika. Volgens deze dimensies heeft een standaardbaan een omtrek van 1609 m, met twee rechte stukken van elk ongeveer 366 m \cite{ofstandard2024}. De resterende lengte van de bochten wordt in formule \ref{eq:booglengte} berekend.

\begin{equation}
    L_{\text{boog}} = \frac{1609 - 2 \times 366}{2} = 439 \text{ m}
    \label{eq:booglengte}
\end{equation}

Beide halve cirkels samen vormen één volledige cirkel met omtrek $2 \times 439 = 878$ m. Door middel van formule \ref{eq:straal} wordt de straal van deze twee bochten berekend.

\begin{equation}
    r = \frac{878}{2\pi} \approx 140 \text{ m}
    \label{eq:straal}
\end{equation}

Bij galop op topsnelheid ($v = 17$ m/s, zoals vermeld in \cite{witte2006effect}) leveren vergelijkingen \ref{eq:hoeksnelheid} en \ref{eq:conversie} de volgende hoeksnelheid op:

\begin{equation}
    \omega_{\text{US}} = \frac{17}{140} \approx 0{,}121 \text{ rad/s} \quad \rightarrow \quad 0{,}121 \times \frac{180}{\pi} \approx 7 \text{ DPS}
    \label{eq:omega_racebaan}
\end{equation}

Een grafische weergave van de baangeometrie is weergegeven in figuur \ref{fig:racebaan}.

\newpage

\begin{figure}[h]
    \centering
    \includegraphics[width=0.85\textwidth]{racetrack_gyro.png}
    \caption{Overzicht van een standaard racebaan voor paarden}
    \label{fig:racebaan}
\end{figure}

De ProRTK tracker zal echter primair ingezet worden bij de \textit{Hong Kong Jockey Club} (HKJC). Hierom wordt bovendien gekeken naar de geometrie van de Sha Tin Racecourse, hun belangrijkste renbaan. De hoofdgrasbaan heeft een omtrek van 1899 m en lange rechte finishstukken van 430 m \cite{wiki_shatin}. Dit resulteert in een bochtstraal van ongeveer $165$ m. Bij een topsnelheid van 17 m/s resulteert dit in een iets lagere hoeksnelheid, doordat de bocht flauwer is dan op de Amerikaanse standaardbaan:

\begin{equation}
    \omega_{\text{Sha Tin}} = \frac{17}{165} \approx 0{,}103 \text{ rad/s} \quad \rightarrow \quad 0{,}103 \times \frac{180}{\pi} \approx 6 \text{ DPS}
    \label{eq:omega_shatin}
\end{equation}

Een top down view van deze renbaan is weergegeven in figuur \ref{fig:sha_tin}

\begin{figure}[h]
    \centering
    \includegraphics[width=0.50\textwidth]{sha_tin.jpg}
    \caption{Top down view van de Sha Tin racebaan \cite{punters_shatin}.}
    \label{fig:sha_tin}
\end{figure}

\newpage

Hoewel de hoeksnelheden op officiële wedstrijd-racebanen met $6$ tot $7$ DPS relatief laag zijn, wordt ook rekening gehouden met onder andere dagelijkse trainingen. Tijdens trainingen ondervinden paarden regelmatig krappere bochten. Om zo dus ruim hierboven te zitten wordt een DPS van minimaal 30 aangenomen. Zo kunnen deze trackers ook eventueel ingezet worden voor andere doeleinden. \\

Een centrale eis is het behoud van de bestaande nauwkeurigheid. In open 
omgevingen levert het huidige RTK-systeem een positienauwkeurigheid van 1 tot 
2 cm bij een volledige RTK Fixed-fix. De toevoeging van de IMU mag deze 
nauwkeurigheid niet negatief beïnvloeden. Tegelijkertijd moet het systeem in staat 
zijn om foutieve GNSS-metingen, veroorzaakt door bijvoorbeeld multipath of 
tijdelijke signaaluitval, te herkennen en af te vangen zodat de 
positieschatting stabiel blijft. Wanneer het GNSS-signaal tijdelijk volledig 
wegvalt, moet de IMU de positieschatting logisch kunnen doorzetten totdat 
het signaal terugkeert. \\

Omdat de metingen live worden bijgehouden, mag er geen merkbare vertraging 
ontstaan. De verwerking van de sensordata moet korter zijn dan het interval 
tussen twee opeenvolgende metingen. De GNSS-ontvanger levert momenteel data 
op 25 Hz, wat neerkomt op een update elke 40 ms. De IMU wordt op een hogere 
frequentie uitgelezen om de positieschatting tussen twee GNSS-metingen bij 
te houden. De exacte frequentie wordt bepaald door de mogelijkheden van de 
hardware en de rekenbelasting op de microcontroller. \\

Om de complexiteit binnen de gestelde projecttijd realistisch te houden, 
beperkt de positiebepaling zich tot het tweedimensionale vlak. 
Hoogtemetingen vallen buiten de scope, aangezien een wedstrijdparcours 
nagenoeg vlak is en hoogtemeting geen onderdeel is van de huidige 
MYLAPS-specificaties. \\

Tabel~\ref{tab:randvoorwaarden} geeft een overzicht van de vastgestelde 
functionele randvoorwaarden van het systeem.

\begin{table}[htbp]
    \centering
    \caption{Functionele randvoorwaarden van het sensorfusie systeem.}
    \label{tab:randvoorwaarden}
    \renewcommand{\arraystretch}{1.3}
    \begin{tabularx}{\textwidth}{|l|X|l|}
        \hline
        \rowcolor{green!20}
        \textbf{Randvoorwaarde} & \textbf{Onderbouwing} & \textbf{Waarde} \\ 
        \hline

        Maximale snelheid &
        Topsnelheid wedstrijdpaard $\approx$ 70,76 km/u~\cite{jockeyclub_racehorses}; 
        10\% veiligheidsmarge aangehouden. &
        $\geq$ 80 km/u \\ \hline

        Accelerometerbereik &
        Piekversnelling op het zwaartepunt bedraagt $\approx$ 2,5 g bij 
        galop~\cite{witte2006effect}. &
        Minimaal 2,5 g \\ \hline

        Gyroscoopbereik &
        Gebaseerd op de verwachte hoeksnelheden bij galop op een 
        trainingsbaan. &
        Minimaal 30 DPS \\ \hline

        GNSS-updatefrequentie &
        Huidige configuratie van de u-blox X20P ontvanger. &
        25 Hz \\ \hline

        Verwerkingstijd &
        Verwerking moet binnen het interval tussen twee GNSS-metingen 
        zijn afgerond om merkbare latency te voorkomen. &
        $\leq$ 40 ms \\ \hline

        Communicatieprotocol &
        Integratie in bestaande MYLAPS-infrastructuur vereist. &
        X2 Link \\ \hline

        Dimensionaliteit &
        Hoogtebepaling valt buiten de scope &
        2D vlak \\ \hline

    \end{tabularx}
\end{table}

Naast de functionele randvoorwaarden zijn de begrippen \textit{robuustheid} en 
\textit{betrouwbaarheid} vertaald naar meetbare prestatiecriteria. Dit is noodzakelijk om objectief te kunnen beoordelen of de positiebepaling daadwerkelijk verbeterd is. In dit project wordt onder \textit{betrouwbaarheid} verstaan dat het systeem onder normale omstandigheden de bestaande nauwkeurigheid behoudt en foutieve metingen niet overneemt. Onder \textit{robuustheid} wordt verstaan dat het systeem bruikbare posities blijft leveren tijdens verstoringen, zoals multipath of een tijdelijke signaaluitval. \\

Voor het bepalen van realistische grenswaarden is gekeken naar wat in de literatuur als haalbaar wordt beschouwd voor vergelijkbare systemen, en naar de wensen vanuit MYLAPS. Onderzoek naar een \textit{tightly-coupled} integratie van RTK met een MEMS-IMU (\textit{Micro-Electro-Mechanical System}) laat zien dat de positiefout tijdens signaaluitval beperkt blijft. Bij een uitval van 15 seconden blijft de nauwkeurigheid op submeterniveau, en bij een uitval van 5 seconden enkele centimeters~\cite{li2017rtk}. Bij een tightly-coupled opzet worden de ruwe GNSS-metingen rechtstreeks met de IMU-data in één filter samengevoegd, wat de hoogste nauwkeurigheid oplevert. Het systeem in dit project gebruikt daarentegen een \textit{loosely-coupled} opzet, aangezien de bewerkte positie en snelheid door de u-blox module worden geleverd. Bovendien is de aanwezige IMU niet geoptimaliseerd voor \textit{dead reckoning} (het uitsluitend op basis van de IMU doorrekenen van de positie wanneer GNSS wegvalt). Een ruimere marge is daardoor realistisch. Het systeem wordt in dit project getoetst op een volledige signaaluitval van 2 seconden, waarvoor in overleg met MYLAPS een grens van één meter is aangenomen. Hoewel de literatuur onder ideale omstandigheden centimeternauwkeurigheid laat zien, wordt deze ruimere grens als voldoende robuust beschouwd. \\

Voor het behoud van de nauwkeurigheid wordt gekeken naar de prestaties van het systeem bij een RTK fix. Onder normale omstandigheden mag de toevoeging van de IMU de huidige nauwkeurigheid van één tot twee centimeter niet verslechteren. \\

Voor het foutdetectiepercentage is in overleg met MYLAPS bepaald dat minimaal 95\% van de grove, foutieve sprongen door multipath moet worden herkend en verworpen. Deze criteria zijn samengevat in tabel~\ref{tab:prestatiecriteria}.

\begin{table}[htbp]
    \centering
    \caption{Prestatiecriteria voor de robuustheid en betrouwbaarheid van het sensorfusie systeem.}
    \label{tab:prestatiecriteria}
    \renewcommand{\arraystretch}{1.3}
    \begin{tabularx}{\textwidth}{|l|X|l|}
        \hline
        \rowcolor{green!20}
        \textbf{Prestatiecriterium} & \textbf{Onderbouwing} & \textbf{Eis} \\ 
        \hline

        Nauwkeurigheidsbehoud (stabiele RTK-fix) &
        In overleg met MYLAPS vastgesteld: de IMU mag de bestaande RTK 
        Fixed-nauwkeurigheid (1--2 cm) niet verslechteren. &
        $\leq$ 2 cm \\ \hline

        Foutdetectiepercentage &
        Foutieve GNSS-metingen door multipath mogen de positieschatting niet 
        meetrekken. Getoetst met softwarematig sprongen van 5 m. &
        $\geq$ 95\% verworpen \\ \hline

        Maximale positiefout bij 2 s uitval &
        Afgeleid van submeterwaarden uit de literatuur~\cite{li2017rtk} en in 
        overleg met MYLAPS. & $\leq$ 1 m \\ \hline

    \end{tabularx}
\end{table}

De randvoorwaarden en prestatiecriteria vormen samen de referentie voor het toetsen van het systeem.

\newpage